\documentclass[12pt,a4paper]{article}
% Estilo compartido para informes. Compilar con LuaLaTeX desde tp1/.
\usepackage[spanish,es-nodecimaldot]{babel}
\usepackage{fontspec}
\setmainfont{Montserrat-Medium.ttf}[
  Path=../assets/fonts/,
  BoldFont=Montserrat-Bold.ttf,
  ItalicFont=Montserrat-MediumItalic.ttf]
\usepackage[a4paper,margin=2.5cm,headheight=16pt]{geometry}
\usepackage{xcolor,graphicx,booktabs,tabularx}
\usepackage{float}
\usepackage{titlesec,fancyhdr,enumitem,caption}
\usepackage{ragged2e}
\usepackage{hyperref}
\definecolor{verde}{HTML}{1D4036}
\definecolor{bronce}{HTML}{D38129}
\definecolor{rosa}{HTML}{D88E77}
\definecolor{calido}{HTML}{EFE4D8}
\definecolor{grafito}{HTML}{1D1D1D}
\hypersetup{colorlinks=true,linkcolor=verde,urlcolor=verde,
  pdftitle={TP1 - Instalación de máquinas virtuales},pdfsubject={Cyberseguridad}}
\AtBeginDocument{\color{grafito}\RaggedRight\fontsize{12}{14.4}\selectfont}
\setlength{\parindent}{0pt}
\setlength{\parskip}{7pt}
\setlist[itemize]{leftmargin=*,itemsep=7pt,topsep=4pt}
\setlist[enumerate]{leftmargin=*,itemsep=6pt}
\titleformat{\section}{\color{verde}\bfseries\fontsize{18}{21.6}\selectfont}{\thesection}{0.65em}{}
\titleformat{\subsection}{\color{verde}\bfseries\large}{\thesubsection}{0.65em}{}
\titleformat{\subsubsection}{\color{verde}\bfseries}{\thesubsubsection}{0.65em}{}
\titlespacing*{\section}{0pt}{16pt}{8pt}
\titlespacing*{\subsection}{0pt}{12pt}{5pt}
\titlespacing*{\subsubsection}{0pt}{10pt}{4pt}
\setcounter{tocdepth}{2}
\pagestyle{fancy}
\fancyhf{}
\fancyhead[L]{\small\color{verde}FCEFyN | UNC}
\fancyhead[R]{\small\color{verde}\materia\ · TP \numeroTP}
\fancyfoot[L]{\scriptsize Informe de laboratorio}
\fancyfoot[R]{\small\thepage}
\renewcommand{\headrulewidth}{0.4pt}
\captionsetup{font=small,labelfont=bf,justification=raggedright,singlelinecheck=false,hypcap=false}
\newcommand{\pendiente}[1]{\par{\small\color{verde}\textbf{Por completar.} #1\par}}
% Si falta una captura se muestra un espacio reservado; no impide compilar.
\newcommand{\captura}[2]{%
  \par\smallskip\begingroup\centering
  \IfFileExists{figuras/#1}{%
    \includegraphics[width=\linewidth,height=5cm,keepaspectratio]{figuras/#1}%
  }{%
    \setlength{\fboxsep}{10pt}%
    \fcolorbox{bronce}{calido}{\parbox[c][1.25cm][c]{\dimexpr\linewidth-2\fboxsep-2\fboxrule\relax}{%
      \centering\small Captura pendiente\par\texttt{\detokenize{#1}}}}%
  }%
  \captionof{figure}{#2}\par\endgroup\smallskip
}

\usepackage{xurl}

\newcommand{\materia}{Cyberseguridad}
\newcommand{\numeroTP}{13}
\newcommand{\tituloTP}{Firma de firmware y Secure Boot v2}
\newcommand{\estudiante}{Brigido Tomas Andres}
\newcommand{\legajo}{38481096}
\newcommand{\carrera}{ICOMP}
\newcommand{\docente}{Javier A. Jorge}


% Metadatos específicos del TP13 y capturas en img/.
\hypersetup{pdftitle={TP13 - Firma de firmware y Secure Boot v2},linkcolor=black,urlcolor=black,citecolor=black}
\renewcommand{\captura}[2]{%
  \begin{figure}[H]
    \centering
    \IfFileExists{img/#1}{%
      \includegraphics[width=0.85\linewidth,height=5cm,keepaspectratio]{img/#1}%
    }{%
      \setlength{\fboxsep}{10pt}%
      \fcolorbox{bronce}{calido}{\parbox[c][1.1cm][c]{0.85\linewidth}{%
        \centering\small Captura pendiente\par\texttt{\detokenize{#1}}}}%
    }%
    \caption{#2}
  \end{figure}
}

\usepackage{caption}
\captionsetup{justification=centering}
\begin{document}
\hypersetup{pageanchor=false}
\begin{titlepage}
  \includegraphics[width=8cm]{../assets/logo-fcefyn-2026.png}
  \vspace{1cm}

  {\color{verde}\small UNIVERSIDAD NACIONAL DE CÓRDOBA\par}
  \vspace{2.1cm}
  {\color{bronce}\bfseries\large TRABAJO PRÁCTICO \numeroTP\par}
  \vspace{0.45cm}
  {\color{verde}\bfseries\fontsize{32}{38.4}\selectfont Firma de firmware\\y Secure Boot v2\par}
  \vspace{0.3cm}
{\small Espressif y SignServer\par}
\vspace{0.35cm}
  {\large \materia\par}
  \vspace{0.8cm}
  {\color{bronce}\rule{2.7cm}{2pt}\par}
  \vfill
  \renewcommand{\arraystretch}{1.45}
  \begin{tabular}{@{}p{3.2cm}p{10cm}@{}}
    Estudiante & \estudiante \\
    Legajo & \legajo \\
    Carrera & \carrera \\
    Docente & \docente \\
  \end{tabular}
  \vfill
  {\small Córdoba, Argentina\par}
\end{titlepage}
\hypersetup{pageanchor=true}

\tableofcontents
\vspace{1cm}
\newpage

\section{Introducción y objetivos}
Esta actividad propone implementar Secure Boot v2 en un dispositivo Espressif e integrar SignServer en el proceso de firma de firmware. Se documentarán la preparación del entorno, la gestión de claves, la firma y las pruebas de arranque.

\section{Marco teórico}
\pendiente{Explicar brevemente qué es Secure Boot v2, qué protege y en qué se diferencia de Flash Encryption. Describir la cadena de confianza y el papel de SignServer.}

\section{Implementación}
\subsection{Entorno de trabajo}
\pendiente{Indicar el modelo y revisión del dispositivo, su compatibilidad con Secure Boot v2, el sistema operativo y las versiones de ESP-IDF y SignServer. Describir su instalación y configuración.}
\captura{01-entorno.png}{Entorno de trabajo y herramientas utilizadas.}

\subsection{Claves y configuración de SignServer}
\pendiente{Documentar la creación del par de claves y la configuración del worker de firma en SignServer.}
\captura{02-signserver.png}{Configuración del worker de firma.}

\newpage
\subsection{Compilación y firma del firmware}
\pendiente{Identificar el ejemplo de ESP-IDF utilizado, como hello\_world, registrar los comandos de compilación y documentar el envío del binario a SignServer y la obtención del firmware firmado.}
\captura{03-firmware-firmado.png}{Compilación y firma del firmware.}

\subsection{Carga y validación en el dispositivo}
\pendiente{Documentar la configuración y activación de Secure Boot v2 y la carga del firmware siguiendo el procedimiento correspondiente al modelo de placa.}
\captura{04-secure-boot.png}{Configuración de Secure Boot v2 en el dispositivo.}

\pendiente{Registrar el arranque con firmware firmado por la clave registrada y contrastarlo con una prueba de firmware no firmado o con firma no válida. Adjuntar los logs y explicar lo observado.}
\captura{05-validacion.png}{Resultados de las pruebas de validación del firmware.}

\section{Problemas encontrados y soluciones}
\pendiente{Indicar los inconvenientes observados, su causa y cómo se resolvieron. Si no hubo inconvenientes, dejarlo registrado.}

\newpage
\section{Cuestionario de validación}
\par\textbf{1. ¿Qué problema de seguridad resuelve el mecanismo de Secure Boot v2?}

\pendiente{Respuesta breve.}

\par\textbf{2. ¿Cuál es la diferencia entre firmar un firmware y cifrarlo?}

\pendiente{Respuesta breve.}

\par\textbf{3. ¿Qué rol cumple SignServer en este proceso?}

\pendiente{Respuesta breve.}

\par\textbf{4. ¿Qué sucede si intentamos cargar un firmware no firmado en un dispositivo con Secure Boot activado?}

\pendiente{Respuesta breve.}

\par\textbf{5. ¿Por qué es importante la gestión centralizada de claves en proyectos de IoT a gran escala?}

\pendiente{Respuesta breve.}

\par\textbf{6. Explica cómo se genera confianza entre el hardware y el firmware en este esquema de seguridad.}

\pendiente{Respuesta breve.}

\par\textbf{7. Menciona un escenario real en el que Secure Boot sería crítico.}

\pendiente{Respuesta breve.}

\section{Reflexión final}
\pendiente{Resumir la importancia de Secure Boot en dispositivos IoT y los aprendizajes de la implementación.}

\section{Referencias}
\begin{itemize}
    \item \textit{Implementación de Secure Boot v2 con Espressif y SignServer}. Consigna local: \texttt{firmadebinarios.md}.
    \item Keyfactor. \href{https://docs.keyfactor.com/signserver/7.0/tutorial-signserver-secure-boot-v2-with-espressif-}{Tutorial: SignServer Secure Boot v2 with Espressif}. Lectura indicada en la consigna.
\end{itemize}
\end{document}
